% PLACEHOLDER introduction — to expand as Phase C/D results land.

\section{Introduction}\label{sec:intro}

[PLACEHOLDER paragraph 1 --- the bias problem.] Classical GWAS, even
with state-of-the-art linear mixed models, omits gene--gene interaction
terms. Yelmen \emph{et~al.}\ \cite{yelmen2026} recently characterised
the resulting bias mathematically: the null $t$-statistic acquires a
mean shift $\mu = \rho\sqrt{\lambda n}/\sqrt{1-\lambda\rho^{2}}$ that
grows with sample size $n$, where $\rho$ is the correlation between
the target SNP and the realised interaction signal and $\lambda$ is
the variance fraction explained by interactions. At biobank scale
($n\geq 100{,}000$) this shift produces 1{,}000s of spurious
genome-wide significant hits even under their strict no-path null
where the target SNP cannot itself participate in the interaction.

[PLACEHOLDER paragraph 2 --- the detection-tool landscape.] Existing
epistasis-detection tools (BOOST \cite{wan2010boost},
MAPIT \cite{crawford2017}, MDR-family approaches) treat the search space
as either an exhaustive pairwise grid (computationally infeasible at
biobank scale) or an unconstrained statistical fishing exercise (low
power after multiple-testing correction). Neither leverages the
biology-typed structure that genome-wide functional resources
(GENCODE, GTEx, STRING, ENCODE, Reactome) make available.

[PLACEHOLDER paragraph 3 --- our contribution.] We present
GraphGWAS-Epistasis, a graph-native toolkit that uses the typed
multi-omics knowledge graph from \cite{estaji2026graphgwas} to
restrict epistasis search to biologically motivated motifs. Four
pairwise detection strategies (M1--M4) operate on this graph-typed
interaction subspace; three higher-order extensions (M5--M7) detect
$k\geq 3$ interactions. The framework also provides a closed-form
$\rho_{\max}$ diagnostic that flags bias-prone target SNPs before
testing.

[PLACEHOLDER paragraph 4 --- empirical headlines, to be filled.]
